\uuid{jM0Y}
\exo7id{4723}
\titre{exo7 4723}
\auteur{quercia}
\organisation{exo7}
\datecreate{2010-03-16}
\isIndication{false}
\isCorrection{true}
\chapitre{Topologie}
\sousChapitre{Topologie de la droite réelle}
\cpge{mpsi-mp2i.limites-continuite.b, pcsi.limites-continuite.b, ptsi.limites-continuite.b, mpsi-mp2i.limites-continuite.c, pcsi.limites-continuite.c, ptsi.limites-continuite.c, mpsi-mp2i.nombres-reels-suites.d, pcsi.nombres-reels-suites.c, ptsi.nombres-reels-suites.c, bcpst1.limites-continuite.c, bcpst1.limites-continuite.d, mp-mpi.topologie-espaces-normes.b, pc.espaces-vectoriels-normes.b, psi.espaces-vectoriels-normes.b, mp-mpi.topologie-espaces-normes.d, pc.espaces-vectoriels-normes.d, psi.espaces-vectoriels-normes.d, pt.fonctions-deux-trois-variables.a, mp-mpi.topologie-espaces-normes.e, pc.espaces-vectoriels-normes.e, psi.espaces-vectoriels-normes.e, pt.fonctions-deux-trois-variables.b}
\module{Analyse}
\niveau{L2}
\difficulte{}

\contenu{
\question{
Soit $f : {[0,1]} \to {[0,1]}$ continue, $u_0 \in {[0,1]}$ et $(u_n)$
la suite des it{\'e}r{\'e}es de $f$ en $u_0$.

On suppose que $u_{n+1} - u_n \xrightarrow[n\to\infty]{} 0$.
Montrer que la suite $(u_n)$ converge vers un point fixe de $f$.
}
\reponse{
L'ensemble des valeurs d'adh{\'e}rence est un intervalle constitu{\'e} de
         points fixes de $f$ $ \Rightarrow $ la suite $(u_n)$ a une seule valeur
         d'adh{\'e}rence.
}
}
